% ============================================================================
%  Inverse Kaluza-Klein Scale Symmetry Framework
%  Document I: An Alternative Solution Framework for the Schrodinger Equation
%  via Scale-Projected Identity Decomposition
%  Compile with pdflatex (twice for refs).
% ============================================================================
\documentclass[11pt, a4paper, oneside]{article}

% --- Geometry ----------------------------------------------------------------
\usepackage[
  top=1.0in, bottom=1.0in, left=1.15in, right=1.15in,
  headheight=14pt
]{geometry}

% --- Fonts / encoding --------------------------------------------------------
\usepackage[T1]{fontenc}
\usepackage{lmodern}
\usepackage{microtype}

% --- Math --------------------------------------------------------------------
\usepackage{amsmath, amssymb, amsthm, mathtools}
\usepackage{bm}
\usepackage{physics}

% --- Tables ------------------------------------------------------------------
\usepackage{booktabs}
\usepackage{array}
\usepackage{tabularx}

% --- Lists -------------------------------------------------------------------
\usepackage{enumitem}
\setlist{nosep, leftmargin=1.5em}

% --- Headers / footers -------------------------------------------------------
\usepackage{fancyhdr}
\pagestyle{fancy}
\fancyhf{}
\fancyhead[L]{\small\textit{Inverse KK Scale Symmetry}}
\fancyhead[R]{\small\textit{Doc I --- Alternative Schrodinger Framework}}
\fancyfoot[C]{\thepage}
\renewcommand{\headrulewidth}{0.4pt}

% --- Section formatting ------------------------------------------------------
\usepackage{titlesec}
\titleformat{\section}
  {\Large\bfseries}
  {\S\thesection}{1em}{}
\titleformat{\subsection}
  {\large\bfseries}
  {\thesubsection}{1em}{}
\titleformat{\subsubsection}
  {\normalsize\bfseries}
  {\thesubsubsection}{1em}{}

% --- Hyperlinks --------------------------------------------------------------
\usepackage[
  colorlinks=true, linkcolor=black, citecolor=black, urlcolor=blue
]{hyperref}

% --- Theorem-like environments -----------------------------------------------
\theoremstyle{definition}
\newtheorem{definition}{Definition}[section]
\newtheorem{axiom}{Axiom}[section]
\theoremstyle{plain}
\newtheorem{proposition}{Proposition}[section]
\newtheorem{theorem}{Theorem}[section]
\theoremstyle{remark}
\newtheorem{remark}{Remark}[section]

% --- Custom operators --------------------------------------------------------
\DeclareMathOperator{\Proj}{Proj}
\newcommand{\Id}{\mathcal{I}}
\newcommand{\IKKRes}{\mathcal{R}}
\newcommand{\Obs}{\mathcal{O}}
\newcommand{\Pack}{\Pi}
\newcommand{\Anch}{\mathcal{A}}
\newcommand{\Scale}{S}

% ============================================================================
\begin{document}

\begin{titlepage}
  \centering
  \vspace*{2cm}
  {\Huge\bfseries Inverse Kaluza-Klein Scale Symmetry\par}
  \vspace{0.8cm}
  {\LARGE Document I\par}
  \vspace{0.5cm}
  {\Large An Alternative Solution Framework for the\\
  Schr\"{o}dinger Equation via\\
  Scale-Projected Identity Decomposition\par}
  \vspace{2cm}
  {\large VSEPR-SIM Theoretical Division\par}
  \vspace{0.4cm}
  {\large Development Day 64\par}
  \vspace{3cm}
  \begin{abstract}
	\noindent
	The standard Schr\"{o}dinger equation operates on a wavefunction
	$\Psi(\mathbf{r}, t)$ defined over classical 3D space plus time.
	This document proposes an alternative framework in which the wavefunction
	is reinterpreted as the \emph{scale-projected shadow} of a fuller
	identity structure $\mathcal{I}$ defined over an extended scale ladder
	$\{S_0, S_1, \ldots, S_D\}$.
	The standard solution is recovered at the $S_3$--$S_4$ projection;
	corrections at higher scales produce measurable residuals that map onto
	dark-sector and quantum-gravity phenomenology.
	The framework does not modify the Schr\"{o}dinger equation's algebra;
	it reinterprets what the wavefunction \emph{is} and why it succeeds and
	fails where it does.
  \end{abstract}
\end{titlepage}

\tableofcontents
\newpage

% ============================================================================
\section{The Standard Framework and Its Projection Assumption}

\subsection{The Schr\"{o}dinger Equation as Usually Stated}

The time-dependent Schr\"{o}dinger equation is:
\begin{equation}
  i\hbar\,\frac{\partial}{\partial t}\Psi(\mathbf{r}, t)
  = \hat{H}\,\Psi(\mathbf{r}, t)
  \label{eq:TDSE}
\end{equation}
where $\hat{H}$ is the Hamiltonian operator, $\mathbf{r} \in \mathbb{R}^3$,
and $\Psi$ is complex-valued. The Born rule assigns probability density:
\begin{equation}
  \rho(\mathbf{r}, t) = |\Psi(\mathbf{r}, t)|^2
\end{equation}

This framework is extraordinarily predictive at the scale $S_3$--$S_4$
(spatial/binding + time-ordered). It fails or requires patching at:
\begin{itemize}
  \item Very high energy / short distance: ultraviolet divergences
  \item Many-body regimes: exponential Hilbert space scaling
  \item Gravitational coupling: no consistent quantum gravity
  \item Cosmological scale: dark matter and dark energy residuals
  \item Measurement: wavefunction collapse has no agreed dynamics
\end{itemize}

\begin{remark}
These failure modes are not random. Under the present framework they are
systematic projection losses: physics that exists at scales beyond $S_4$
but is not captured by the $\Psi(\mathbf{r},t)$ representation.
\end{remark}

\subsection{The Implicit Projection Assumption}

The standard equation contains a hidden assumption:
\begin{equation}
  \boxed{
	\Psi(\mathbf{r},t) \equiv P_{S_3 \cup S_4}(\mathcal{I})
  }
  \label{eq:implicit-projection}
\end{equation}
That is, the wavefunction \emph{is} the identity state. This is the
assumption the present framework relaxes. Instead:
\begin{equation}
  \Psi(\mathbf{r},t) = P_{S_3 \cup S_4}(\mathcal{I})
  \quad \text{with} \quad
  \mathcal{I} \supset P_{S_3 \cup S_4}(\mathcal{I})
  \label{eq:relaxed}
\end{equation}
The residual is:
\begin{equation}
  \mathcal{R}_{S_3} = \mathcal{U}_{S_3}(\mathcal{I}) - \mathcal{E}_{S_3}(\Psi)
  \label{eq:residual}
\end{equation}
where $\mathcal{U}_{S_3}$ is the full unprojected identity content at scale
$S_3$ and $\mathcal{E}_{S_3}$ is the classical expectation recovery.

% ============================================================================
\section{Identity Structure and Scale Ladder}

\subsection{Scale Ladder Definition}

\begin{definition}[Scale entry]
  Each scale level $S_n$ is a triple:
  \begin{equation}
	S_n = (D_n,\; M_n,\; L_n)
  \end{equation}
  where $D_n$ is the degree of freedom / dimensional permission added at
  scale $n$, $M_n$ is the associated mediator or force carrier, and $L_n$
  is the observational limit (the scale below which $S_n$ becomes relevant).
\end{definition}

\begin{table}[h]
  \centering
  \begin{tabularx}{\linewidth}{c l l l}
	\toprule
	Scale & Meaning & Mediator & Observational limit \\
	\midrule
	$S_0$ & Binary existence $E \in \{0,1\}$ & --- & Planck? \\
	$S_1$ & Linear vector permission & --- & Sub-nuclear \\
	$S_2$ & Relational / angular & Gluons (color) & Nuclear \\
	$S_3$ & Spatial / binding structure & Photons (EM) & Atomic \\
	$S_4$ & Time-ordered transformation & $W, Z$ bosons & Weak \\
	$S_5$ & Mass-energy curvature & Graviton (hypothetical) & Cosmological \\
	$S_D$ & Dark-scale permission $w$ & $X_D$ (unknown) & $L_\text{EM}$ \\
	\bottomrule
  \end{tabularx}
  \caption{Scale ladder. Time $t$ is the sequencing parameter of
  $\mathcal{I}(t)$, not a scale level itself.}
  \label{tab:scale-ladder}
\end{table}

\subsection{Identity Bundle and Identity Anchor}

A particle is not a point. It is an identity bundle:
\begin{equation}
  \mathcal{I}_i(t) = \{c_{i1}, c_{i2}, \ldots, c_{im}\}
  \label{eq:identity-bundle}
\end{equation}
where each component $c_{ik}$ carries a scale-level tag.

The identity anchor (generalized center) is:
\begin{equation}
  \mathcal{A}_i(t)
  = \frac{\sum_k \alpha_{ik}\, c_{ik}(t)}{\sum_k \alpha_{ik}}
  = \frac{\bm{\alpha}_i^\top \mathbf{C}_i}{\bm{\alpha}_i^\top \mathbf{1}}
  \label{eq:anchor}
\end{equation}

The observed 3D position is only the projection:
\begin{equation}
  \vec{r}_i^{\,\text{obs}} = P_{3D}\!\left(\mathcal{A}_i\right)
\end{equation}

The fuzz radius measures internal spread at the observed scale:
\begin{equation}
  r_f^2 =
  \frac{\sum_k \alpha_k \left\|P_{3D}(c_k) - \vec{r}^{\,\text{obs}}\right\|^2}
	   {\sum_k \alpha_k}
  \label{eq:fuzz-radius}
\end{equation}

\begin{definition}[Fuzz ratio]
  \begin{equation}
	\Phi_S = \frac{r_f}{\ell_S}
  \end{equation}
  where $\ell_S$ is the characteristic length of scale $S$.
  \begin{align*}
	\Phi_S \ll 1 &\implies \text{point-like at scale } S \\
	\Phi_S \sim 1 &\implies \text{fuzz-zone} \\
	\Phi_S > 1  &\implies \text{non-point-defined at scale } S
  \end{align*}
\end{definition}

% ============================================================================
\section{The Alternative Wavefunction}

\subsection{Extended State Space}

In the standard framework, the state space is:
\begin{equation}
  \Psi: \mathbb{R}^3 \times \mathbb{R} \to \mathbb{C}
\end{equation}

In the extended framework, the full identity state lives in:
\begin{equation}
  \Psi_\text{full}: \mathbb{R}^3 \times w \times \mathbb{R} \to \mathbb{C}
  \qquad
  \Psi_\text{full} = \Psi(x, y, z, w;\, t)
  \label{eq:extended-state}
\end{equation}
where $w$ is the dark-scale coordinate. The observed wavefunction is the
projection:
\begin{equation}
  \Psi_\text{obs}(\mathbf{r}, t)
  = \int \Psi_\text{full}(\mathbf{r}, w, t)\, K_\text{EM}(w)\, dw
  \label{eq:projected-wavefunction}
\end{equation}
with $K_\text{EM}(w)$ the electromagnetic projection kernel.

\subsection{The Extended Schr\"{o}dinger Equation}

Replace $\hat{H}$ with the extended Hamiltonian:
\begin{equation}
  \hat{H}_\text{eff} = \hat{H}_{S_3} + \hat{H}_{S_4} + \hat{H}_w
  \label{eq:extended-H}
\end{equation}
where $\hat{H}_w$ is the dark-scale Hamiltonian term. The extended equation
becomes:
\begin{equation}
  \boxed{
	i\hbar\,\frac{\partial}{\partial t}\Psi_\text{full}(\mathbf{r}, w, t)
	= \hat{H}_\text{eff}\,\Psi_\text{full}(\mathbf{r}, w, t)
  }
  \label{eq:extended-TDSE}
\end{equation}

Integrating out $w$ with kernel $K_\text{EM}(w)$ recovers~\eqref{eq:TDSE}
exactly when $\hat{H}_w$ is negligible, plus a correction term:
\begin{equation}
  i\hbar\,\frac{\partial \Psi_\text{obs}}{\partial t}
  = \hat{H}_{S_3}\,\Psi_\text{obs}
  + \underbrace{\int \hat{H}_w\,\Psi_\text{full}\, K_\text{EM}(w)\, dw}
			   _{\displaystyle\delta\hat{H}_w[\Psi_\text{obs}]}
  \label{eq:corrected-TDSE}
\end{equation}

\begin{remark}
  $\delta\hat{H}_w$ is the residual Hamiltonian term invisible to EM
  projection. It encodes exactly the physics that the standard equation
  misses. In regimes where $\hat{H}_w \approx 0$, standard QM is
  recovered without modification.
\end{remark}

\subsection{Cancellation and Hidden Active Structure}

A critical property of the extended framework is that observable channel
cancellation does not imply internal inactivity.

Define the projected observable $I_j$ and hidden intensity $H_j$ for
component $j$:
\begin{align}
  I_j &= \bm{\alpha}^\top \mathbf{c}_j \label{eq:observable-channel}\\
  H_j &= \mathbf{c}_j^\top \mathbf{W}\, \mathbf{c}_j \label{eq:hidden-intensity}
\end{align}

The hidden-active condition:
\begin{equation}
  \boxed{I_j = 0, \quad H_j > 0}
  \label{eq:hidden-active}
\end{equation}

\textbf{Example: neutron-like charge cancellation.}
\begin{equation}
  I_\text{charge} = +\tfrac{2}{3} - \tfrac{1}{3} - \tfrac{1}{3} = 0
\end{equation}
\begin{equation}
  H_\text{charge}
  = \left(\tfrac{2}{3}\right)^2
  + \left(-\tfrac{1}{3}\right)^2
  + \left(-\tfrac{1}{3}\right)^2
  = \tfrac{2}{3} \neq 0
\end{equation}
The particle appears charge-neutral (zero observable channel) but has
nonzero internal charge intensity. This structure is hidden from the
$K_\text{EM}$ projection but present in $\hat{H}_w$.

This generalizes to:
\begin{center}
  \begin{tabular}{ll}
	\toprule
	Observable channel & Hidden-active mechanism \\
	\midrule
	Electric charge & Quark charge cancellation \\
	Color charge & Color neutrality (confinement) \\
	Spin & Singlet pairing \\
	Linear momentum & Bound-state internal motion \\
	Angular momentum & $L + S = J = 0$ nuclear ground states \\
	Magnetic moment & Quenched orbital contributions \\
	Defect pairs & Frenkel / Schottky charge balance \\
	Mechanical stress & Residual stress in balanced structures \\
	\bottomrule
  \end{tabular}
\end{center}

% ============================================================================
\section{Projection Failure and Exactness}

\subsection{Many-to-One Projection}

\begin{definition}[Projection non-injectivity]
  A projection $P_S$ is non-injective when:
  \begin{equation}
	P_S(\mathcal{I}_a) = P_S(\mathcal{I}_b)
	\quad \text{while} \quad
	\mathcal{I}_a \neq \mathcal{I}_b
  \end{equation}
\end{definition}

\begin{theorem}[Exactness failure]
  When $P_S$ is non-injective:
  \begin{equation}
	P_S^{-1}(\Obs_S) \text{ is not unique.}
  \end{equation}
  Exact reconstruction from observation at scale $S$ is impossible.
\end{theorem}

\begin{proof}[Sketch]
  Two distinct identity states $\mathcal{I}_a \neq \mathcal{I}_b$ produce
  identical observations $\Obs_S$. No measurement at scale $S$ can
  distinguish them. A unique inverse map would require information not
  present at scale $S$. \qed
\end{proof}

\subsection{Destructive Inverse Packing}

Packing is the forward process:
\begin{equation}
  \mathbf{C}^{(n)} \xrightarrow{\Pi_n} \mathbf{I}^{(n+1)}
  \label{eq:packing}
\end{equation}
The inverse (annihilation / separation):
\begin{equation}
  \mathbf{I}^{(n+1)} \xrightarrow{\mathcal{A}} \mathbf{C}_\text{out}^{(n)}
  \label{eq:annihilation}
\end{equation}
necessarily \emph{destroys or transforms} the original packed identity.

\begin{proposition}
  Complete hidden-variable maps require destructive inverse packing.
  Therefore exact mapping cannot preserve the system being mapped.
\end{proposition}

This connects to the quantum no-cloning theorem and to the irreversibility
of measurement. It is not an epistemological limitation; it is a structural
consequence of packing.

% ============================================================================
\section{Entanglement as a Relation-Particle}

\subsection{The Relation Object}

Two entangled particles $A$ and $B$ share a relation object:
\begin{equation}
  \mathcal{R}_{AB} = \bigl[0,\, 0,\, 0,\, 0\;\big|\;\mathbf{T}_{AB}\bigr]
  \label{eq:relation-object}
\end{equation}
The first four entries are spacetime displacement channels:
\begin{equation}
  \Delta x = \Delta y = \Delta z = \Delta t = 0
\end{equation}
The trinary relation vector:
\begin{equation}
  \mathbf{T}_{AB} \in \{-1, 0, +1\}^m
\end{equation}
carries the real correlational content.

\begin{remark}
  $[0,0,0,0]$ means \emph{no spacetime displacement}, not
  \emph{no physical existence}. The relation object is real,
  deterministic, and non-spatiotemporal.
\end{remark}

\subsection{Pair Decomposition and Measurement}

The joint identity decomposes as:
\begin{equation}
  \mathcal{I}_{AB} \;\to\; \mathcal{I}_A + \mathcal{I}_B + \mathcal{R}_{AB}
  \label{eq:pair-decomposition}
\end{equation}

Measurement outcomes are scale-projected with the relation object:
\begin{align}
  O_A &= P_a\!\left(\mathcal{I}_A,\, \mathcal{R}_{AB}\right) \\
  O_B &= P_b\!\left(\mathcal{I}_B,\, \mathcal{R}_{AB}\right)
\end{align}

The correlation function:
\begin{equation}
  E(a,b) = \langle O_A\, O_B \rangle
\end{equation}
reproduces quantum-mechanical predictions because $\mathcal{R}_{AB}$ is
a fixed structure in identity-space, not a propagating signal.

% ============================================================================
\section{Dark Matter as Higher-Scale Projection Residual}

\subsection{The Projection Gap}

\begin{equation}
  \rho_\text{dark}(x,y,z;\,t)
  = P_G(\rho) - P_\text{EM}(\rho)
  \label{eq:dark-def}
\end{equation}

In kernel form, integrating over the dark coordinate $w$:
\begin{equation}
  \boxed{
	\rho_\text{dark}(x,y,z;\,t)
	= \int \rho(x,y,z,w;\,t)
	  \left[K_G(w) - K_\text{EM}(w)\right] dw
  }
  \label{eq:dark-kernel}
\end{equation}

$K_G(w)$ is the gravitational projection kernel (couples to all scales);
$K_\text{EM}(w)$ is the electromagnetic projection kernel (truncates at
$L_\text{EM}$).

\subsection{Dark-Scale Entry}

The dark-scale level in the ladder:
\begin{equation}
  S_D = (w,\; X_D,\; L_\text{EM})
\end{equation}
where $w$ is the dark-scale degree of freedom, $X_D$ is the (presently
unknown) dark mediator, and $L_\text{EM}$ marks where electromagnetic
projection fails.

\textbf{Central claim:}
\begin{equation}
  \boxed{
	\text{dark matter} = \text{gravitational signature of
	an unresolved higher-scale permission}
  }
\end{equation}

This is distinct from standard extra-dimensional models:
\begin{align*}
  \text{Standard:} &\quad \text{dimension first} \to \text{dark behavior} \\
  \text{Inverse KK:} &\quad \text{force-scale residuals first}
							  \to \text{inferred dimension}
\end{align*}

% ============================================================================
\section{Effective Proper Time with Dark-Scale Correction}

\subsection{Classical and Quantum Proper Time}

Classical relativistic proper time:
\begin{equation}
  d\tau_\text{4D}
  = dt\sqrt{1 - \frac{v^2}{c^2} + \frac{2\Phi_G}{c^2}}
  \label{eq:classical-proper-time}
\end{equation}

Recent optical ion clock work establishes:
\begin{equation}
  \hat{\tau} \equiv \tau(\hat{x}, \hat{p})
  \label{eq:quantum-proper-time}
\end{equation}
Proper time becomes quantum-linked; $\hat{\tau}$ is an observable with
spread $R_\tau = \hat{\tau} - \langle\hat{\tau}\rangle$.

\subsection{Dark-Scale Extension}

Extend to:
\begin{equation}
  \boxed{
	d\tau_\text{eff}
	= dt\sqrt{
		1 - \frac{v^2}{c^2}
		+ \frac{2\Phi_G}{c^2}
		+ \chi_w(w)
	  }
  }
  \label{eq:effective-proper-time}
\end{equation}
where $\chi_w(w)$ is the dark-scale metric correction term.

In differential form:
\begin{equation}
  d\tau_\text{eff} = d\tau_\text{4D} + d\tau_w
  \label{eq:tau-decomp}
\end{equation}

The full proper-time identity structure:
\begin{equation}
  \mathcal{I}_\tau
  = \left(\hat{x},\, \hat{p},\, \hat{\tau},\, H_c,\, H_m,\, w\right)
  \label{eq:proper-time-identity}
\end{equation}
where $H_c$ is the clock Hamiltonian and $H_m$ is the motional Hamiltonian.
The clock-motion relation object:
\begin{equation}
  \mathcal{R}_{cm} = [0, 0, 0, 0 \mid \mathbf{T}_{cm}]
\end{equation}

\begin{remark}
  $\chi_w(w)$ remains formally undefined here. Constraining it requires
  either: (a) a specific model for $\hat{H}_w$, or (b) comparison with
  clock-discrepancy data in strong gravitational fields or extreme
  velocities where 4D relativity is insufficient.
\end{remark}

% ============================================================================
\section{Recovery of Standard Quantum Mechanics}

The critical consistency check: the standard Schr\"{o}dinger equation must
be exactly recovered in the appropriate limit.

\begin{proposition}[Standard QM recovery]
  When $\hat{H}_w \approx 0$ (dark-scale Hamiltonian negligible) and
  $\chi_w(w) \approx 0$ (proper-time correction negligible), the
  extended framework reduces exactly to:
  \begin{align}
	\Psi_\text{obs}(\mathbf{r},t) &\approx \Psi(\mathbf{r},t) \\
	i\hbar\,\partial_t\Psi &= \hat{H}_{S_3}\Psi
  \end{align}
  which is~\eqref{eq:TDSE}.
\end{proposition}

The framework is therefore a \emph{conservative extension}. No existing
prediction is contradicted; corrections appear only where the extended
terms are non-negligible: strong gravity, cosmological scales, and
dark-sector phenomenology.

% ============================================================================
\section{Open Problems}

\begin{enumerate}
  \item \textbf{Form of $K_\text{EM}(w)$ and $K_G(w)$.}
	These projection kernels determine the magnitude of all dark-sector
	corrections. They are presently unconstrained.

  \item \textbf{Form of $\hat{H}_w$.}
	Without a specific dark-scale Hamiltonian the theory has no
	new quantitative predictions.

  \item \textbf{Form of $\chi_w(w)$.}
	The dark-scale metric correction in~\eqref{eq:effective-proper-time}
	needs a physical derivation or at minimum a phenomenological ansatz.

  \item \textbf{Trinary relation vector $\mathbf{T}_{AB}$.}
	The space $\{-1, 0, +1\}^m$ must be given a dynamics. How does
	$\mathbf{T}_{AB}$ evolve under local operations?

  \item \textbf{Packing operator $\Pi_n$.}
	The map from $\mathbf{C}^{(n)}$ to $\mathbf{I}^{(n+1)}$ is stated
	as an axiom. Its explicit form at each scale level is needed.

  \item \textbf{Connection to existing field theory.}
	The extended equation~\eqref{eq:extended-TDSE} must be
	Lorentz-covariant. The dark coordinate $w$ must respect causality
	or explain why it does not.
\end{enumerate}

% ============================================================================
\appendix
\section{Notation Summary}

\begin{tabular}{ll}
  \toprule
  Symbol & Meaning \\
  \midrule
  $\mathcal{I}$ & Full identity state \\
  $P_S$ & Projection onto scale $S$ \\
  $\mathcal{R}_S$ & Residual at scale $S$: $\mathcal{U}_S(\mathcal{I}) - \mathcal{E}_S(\Obs_S)$ \\
  $S_n = (D_n, M_n, L_n)$ & Scale entry: DOF, mediator, observational limit \\
  $w$ & Dark-scale degree of freedom \\
  $\mathcal{A}_i$ & Identity anchor (eq.~\ref{eq:anchor}) \\
  $r_f$ & Fuzz radius (eq.~\ref{eq:fuzz-radius}) \\
  $\Phi_S$ & Fuzz ratio \\
  $\Pi_n$ & Packing operator \\
  $\mathcal{R}_{AB}$ & Relation-particle object (eq.~\ref{eq:relation-object}) \\
  $\mathbf{T}_{AB}$ & Trinary relation vector \\
  $K_G(w), K_\text{EM}(w)$ & Gravitational / EM projection kernels \\
  $\chi_w(w)$ & Dark-scale metric correction \\
  $d\tau_\text{eff}$ & Effective proper time (eq.~\ref{eq:effective-proper-time}) \\
  \bottomrule
\end{tabular}

\end{document}
